%%%%%%%%%%%%%%%%%%%%%%%%%%%%%%%%%%%%%%%%%%%%%%%%%%%%%%%%%%%%%%%%%%%%%%%%%%%%%%%
%% Informe del proyecto soc-purple-team-wazuh
%% Hecho con la plantilla personal mapc-proyecto.cls
%% Compilar: pdflatex -> biber -> makeglossaries -> pdflatex x2
%% Las capturas se leen de ../evidencia/ (la carpeta del repo).
%%%%%%%%%%%%%%%%%%%%%%%%%%%%%%%%%%%%%%%%%%%%%%%%%%%%%%%%%%%%%%%%%%%%%%%%%%%%%%%
\documentclass{mapc-proyecto}
\graphicspath{{../evidencia/}}
% \evidenciaH: igual que \evidencia, pero fija la captura en su sitio
\newcommand{\evidenciaH}[4][0.92\textwidth]{%
  \begin{figure}[H]\centering
  {\setlength{\fboxsep}{0pt}\setlength{\fboxrule}{0.6pt}%
   \color{mapclinea}\fbox{\includegraphics[width=#1]{#2}}}%
  \caption{#3}\label{#4}
  \end{figure}}

% ---------------------------- METADATOS --------------------------------------
\proyecto{SOC Purple Team en vivo con Wazuh}
\subtitulo{Un SOC real en AWS, atacado desde Kali, con un panel que muestra en vivo
la cobertura de MITRE ATT\&CK sobre alertas reales}
\tipoproyecto{Proyecto personal \textbullet\ Ciberseguridad defensiva}
\repositorio{\url{https://github.com/AnyelokDev/soc-purple-team-wazuh}}
\stack{Wazuh 4.14.7, Amazon EC2 (Amazon Linux 2023), Python 3 (librería estándar),
HTML/JS, Kali Linux, Fable 5.1 vía Claude Code}
\contexto{Construido y presentado en el evento Fable 5.1 Build Day (septiembre de 2026)}
\estadoproyecto{Terminado}
\versiondoc{1.0}
\clasificacion{Público \textbullet\ portafolio}
\cabecera{SOC Purple Team en vivo con Wazuh}

\addbibresource{referencias.bib}

% ------------------------------ SIGLAS ---------------------------------------
\newacronym{soc}{SOC}{Centro de Operaciones de Seguridad}
\newacronym{siem}{SIEM}{gestión de información y eventos de seguridad}
\newacronym{sse}{SSE}{Server-Sent Events}
\newacronym{ec2}{EC2}{Amazon Elastic Compute Cloud}
\newacronym{ttp}{TTP}{tácticas, técnicas y procedimientos}

\begin{document}
\portada

%%%%%%%%%%%%%%%%%%%%%%%%%%%%%%%%%%%%%%%%%%%%%%%%%%%%%%%%%%%%%%%%%%%%%%%%%%%%%%%
%% CAPA 1 — RESUMEN
%%%%%%%%%%%%%%%%%%%%%%%%%%%%%%%%%%%%%%%%%%%%%%%%%%%%%%%%%%%%%%%%%%%%%%%%%%%%%%%
\begin{resumen}
Monté un \gls{soc} pequeño pero real: Wazuh 4.14.7 en una instancia de Amazon
\gls{ec2}, atacado desde una máquina Kali, y un panel web que lee las alertas de
Wazuh en vivo y pinta sobre ellas una cuadrícula de MITRE ATT\&CK con un porcentaje
de cobertura. La idea era mostrar en una demo de pocos minutos el ciclo completo
de un equipo \emph{purple}: se ataca, se mira qué detecta el SOC, se identifica
lo que falta y se cierra con una regla nueva.

El set de la demo son siete técnicas. Cinco las cubrieron las reglas de fábrica de
Wazuh en cuanto hubo actividad. Las otras dos corresponden al \textbf{escaneo de
red con \texttt{nmap}} (T1595 y T1046). El escaneo funcionó: llegó al servidor,
identificó los servicios expuestos y dejó sus huellas en el log de SSH. Lo que no
existía era una regla que interpretara esos eventos, así que para el Wazuh de
fábrica el escaneo era invisible. Con un decoder y tres reglas de correlación
nuevas, el escaneo pasó a detectarse y el panel subió del 43\,\% al 71\,\%, hasta
llegar al \textbf{100\,\% de cobertura con alertas reales}.

\textbf{El aporte principal del proyecto es cómo se escribieron esas reglas.} Yo
diseñé el flujo, monté el SOC, lancé los ataques y dirigí el trabajo. Dentro de ese
flujo, Fable 5.1 (vía Claude Code) generó de forma autónoma el decoder y las reglas
de detección, y las fue mejorando a partir de lo que el Wazuh real registraba hasta
que el escaneo quedó detectado. También escribió la regla de persistencia (T1136)
y el panel. Todo se validó contra alertas reales. Durante la demo el panel no
consulta al modelo; todo lo que muestra sale de \texttt{alerts.json}.

Queda pendiente un modo de anonimización en el panel, para poder grabarlo sin
que aparezcan IPs ni nombres de host reales.
\end{resumen}

\seccionsn{Qué demuestra este proyecto}

\begin{table}[!htb]\centering\small
\begin{tabularx}{\textwidth}{L{4.3cm} Y}
\toprule
\filacab \thc{Competencia} & \thc{Dónde se ve} \\
\midrule
Desplegar y endurecer un SIEM en la nube & Wazuh todo en uno en EC2, SSH solo con clave (\cref{sec:arq}). \\
Detección como código con un modelo de IA & Fable 5.1 generó de forma autónoma el decoder \texttt{sshd-banner-invalid} y las reglas 100098--100100 dentro del flujo que diseñé (\cref{sec:reglas}). \\
Mapeo a MITRE ATT\&CK & Matriz de cobertura de las siete técnicas (\ref{anx:matriz}). \\
Análisis con evidencia & El caso T1136, donde la evidencia contradijo la hipótesis inicial (\cref{sec:problemas}). \\
Desarrollo de herramientas propias & Panel en vivo con Python y SSE, sin dependencias (\cref{sec:panel}). \\
OPSEC al publicar & Capturas censuradas y revisadas; tres fugas que el OCR no vio (\cref{sec:seguridad}). \\
\bottomrule
\end{tabularx}
\caption{Competencias y dónde se respaldan.}\label{tab:competencias}
\end{table}

\seccionsn{Resultados clave}

\begin{hallazgo}[Resultado 1 --- 100\,\% de cobertura del set con alertas reales]
Las siete técnicas del set (T1595, T1046, T1110, T1078, T1021, T1548 y T1136)
quedaron cubiertas por alertas reales de Wazuh.\\[3pt]
\textbf{Evidencia:} capturas 14, 18 y 22 (\ref{anx:evidencia}).
\end{hallazgo}

\begin{hallazgo}[Resultado 2 --- el escaneo pasó de invisible para Wazuh a detectado]
El escaneo con \texttt{nmap} se completó contra el servidor (captura 05), pero el
Wazuh de fábrica no tenía una regla que interpretara esos eventos, así que no
generaba alerta. Con el decoder y la regla de correlación 100100, escritos por
Fable, el mismo escaneo produce una alerta de nivel 8 mapeada a T1046 y T1595, y
la cobertura del panel sube del 43\,\% al 71\,\%.\\[3pt]
\textbf{Evidencia:} capturas 05 a 08.
\end{hallazgo}

\begin{hallazgo}[Resultado 3 --- publicación sin datos reales]
El repositorio no contiene IPs, nombres de host ni credenciales reales. Las
capturas usan direcciones de documentación (RFC 5737), y en la revisión final
aparecieron tres fugas que la verificación con OCR había pasado por alto.\\[3pt]
\textbf{Evidencia:} \cref{sec:seguridad}.
\end{hallazgo}

\newpage
\tableofcontents
\newpage

%%%%%%%%%%%%%%%%%%%%%%%%%%%%%%%%%%%%%%%%%%%%%%%%%%%%%%%%%%%%%%%%%%%%%%%%%%%%%%%
%% CAPA 2 — TÉCNICA
%%%%%%%%%%%%%%%%%%%%%%%%%%%%%%%%%%%%%%%%%%%%%%%%%%%%%%%%%%%%%%%%%%%%%%%%%%%%%%%
\section{Contexto y objetivo}

El proyecto nació para el \emph{Fable 5.1 Build Day}, un evento en el que había
que construir algo con el modelo Fable 5.1 y presentarlo en vivo. Elegí un tema
de mi área: automatizar parte del trabajo de un analista de \gls{soc}, que es
revisar qué técnicas de ataque detecta el \gls{siem}, cuáles no, y escribir las
reglas que faltan.

El objetivo concreto fue una demo en la que el público viera:

\begin{enumerate}
  \item Ataques reales contra un servidor real, no un video ni datos simulados.
  \item Cada alerta de Wazuh apareciendo en pantalla en el momento en que ocurre,
        con su técnica de MITRE ATT\&CK \parencite{mitre-attack}.
  \item Un porcentaje de cobertura sobre un set cerrado de técnicas, que sube a
        medida que se detectan.
  \item Un hueco de detección y su cierre con una regla escrita para el caso.
\end{enumerate}

\subsection{Reparto del trabajo}

Como el proyecto se hizo con un modelo de IA, dejo claro qué hizo cada parte:

\begin{table}[!htb]\centering\small
\begin{tabularx}{\textwidth}{L{3.4cm} Y}
\toprule
\filacab \thc{Quién} & \thc{Qué hizo} \\
\midrule
Yo & Montaje de la infraestructura (EC2, Wazuh, endurecimiento de SSH), definición del set de técnicas y de la demo, dirección y validación del trabajo del modelo, ejecución de los ataques, revisión OPSEC y presentación. \\
Fable 5.1 (vía Claude Code) & Escribió el panel (\texttt{server.py} e \texttt{index.html}), el decoder y las reglas de escaneo, y la regla 100200, que detecta la creación de cuentas locales, la mapea a T1136 (persistencia) y la eleva a nivel 10. \\
\bottomrule
\end{tabularx}
\caption{Reparto del trabajo.}\label{tab:reparto}
\end{table}

\section{Arquitectura}\label{sec:arq}

Todo el lado defensivo vive en una sola instancia \gls{ec2}. Wazuh corre en modo
todo en uno (manager, indexer y dashboard) y vigila el propio servidor: los
eventos de SSH, PAM y del sistema del host son la telemetría. El panel en vivo
corre en la misma máquina y solo lee el archivo de alertas del manager.

\begin{figure}[H]\centering
\begin{tikzpicture}[node distance=0.55cm and 0.7cm, font=\sffamily\small]
  \node[nodooscuro, minimum width=2.3cm] (kali) {Kali Linux\\[-1pt]{\scriptsize\mdseries atacante}};
  \node[nodo, right=1.2cm of kali, minimum width=2.6cm] (host) {Sistema del host\\[-1pt]{\scriptsize sshd, PAM, sudo}};
  \node[nodo, right=of host, minimum width=2.8cm] (wazuh) {Wazuh manager\\[-1pt]{\scriptsize decoders + reglas}};
  \node[nodo, below=of wazuh, minimum width=2.8cm] (json) {\texttt{alerts.json}};
  \node[nodoacento, left=of json, minimum width=2.6cm] (srv) {\texttt{server.py}\\[-1pt]{\scriptsize tail -F + SSE}};
  \node[nodooscuro, below=0.9cm of kali, minimum width=2.3cm] (nav) {Navegador\\[-1pt]{\scriptsize\mdseries index.html}};
  \begin{scope}[on background layer]
    \node[zona, fit=(host)(wazuh)(json)(srv), label={[font=\sffamily\scriptsize\bfseries,text=mapcmorado]above:EC2 \textbullet\ Amazon Linux 2023}] {};
  \end{scope}
  \draw[flecha] (kali) -- node[etiq,above] {ataques} (host);
  \draw[flecha] (host) -- node[etiq,above] {logs} (wazuh);
  \draw[flecha] (wazuh) -- node[etiq,right] {alerta} (json);
  \draw[flecha] (json) -- (srv);
  \draw[flechaacento] (srv) -- node[etiq,above] {SSE :8080} (nav);
\end{tikzpicture}
\caption{Arquitectura del laboratorio.}\label{fig:arq}
\end{figure}

\begin{table}[H]\centering\small
\begin{tabularx}{\textwidth}{L{3.3cm} Y}
\toprule
\filacab \thc{Componente} & \thc{Detalle} \\
\midrule
Servidor & EC2 \texttt{t3.large} con Amazon Linux 2023 e IP fija (Elastic IP). SSH solo con clave; usuario \texttt{ec2-user}. \\
SIEM & Wazuh 4.14.7 todo en uno, instalado con el asistente oficial \parencite{wazuh-quickstart}. \\
Panel & \texttt{dashboard/server.py} (Python, solo librería estándar) e \texttt{index.html}. Escucha en el puerto 8080. \\
Reglas propias & \texttt{wazuh-rules/}: un decoder, tres reglas de escaneo y la regla 100200. \\
Atacante & Kali Linux en VirtualBox. El escaneo y los intentos de acceso salen de ahí; la creación de cuentas se hizo en la propia EC2. \\
\bottomrule
\end{tabularx}
\caption{Componentes.}\label{tab:componentes}
\end{table}

\section{Implementación}

\subsection{Reglas y decoder propios}\label{sec:reglas}

Wazuh ya trae reglas para casi todo el set. Lo que faltaba era una forma de
detectar el escaneo de red, que sí llegaba al servidor.
El problema tenía dos partes (\cref{sec:problemas}): la regla de fábrica para
<<conexión cerrada>> tiene nivel 0, así que no genera alerta, y las sondas no-SSH
de \texttt{nmap -A} quedan en el log de OpenSSH 9.x como
\texttt{banner exchange: \ldots\ invalid format}, un mensaje que ningún decoder de
fábrica interpreta, así que tampoco se extrae la IP de origen.

La solución tiene tres piezas, siguiendo la sintaxis de reglas de Wazuh
\parencite{wazuh-rules}:

\begin{table}[!htb]\centering\small
\begin{tabularx}{\textwidth}{C{1.6cm} C{1.1cm} Y L{2.4cm}}
\toprule
\filacab \thc{ID} & \thc{Nivel} & \thc{Qué hace} & \thc{MITRE} \\
\midrule
decoder & --- & \texttt{sshd-banner-invalid}: extrae \texttt{srcip} y \texttt{srcport} del mensaje \texttt{banner exchange}. & --- \\
100099 & 1 & Hija de la 5722 (<<conexión cerrada>>). La sube a nivel 1 para que cuente como evento base. & --- \\
100098 & 1 & Sonda con protocolo inválido en el puerto SSH. Evento base. & --- \\
100100 & 8 & Dos o más eventos base desde la misma IP en 5 minutos: escaneo. & \chip{T1046} \chip{T1595} \\
100200 & 10 & Hija de la 5902 (<<cuenta creada>>): detecta la creación de cuentas locales como persistencia, la mapea a T1136 y la sube de nivel 8 a 10. & \chip{T1136} \\
\bottomrule
\end{tabularx}
\caption{Decoder y reglas propias.}\label{tab:reglas}
\end{table}

Las otras técnicas las cubre Wazuh de fábrica \parencite{wazuh-sshd-rules}:
5710 y 5712 (usuario inexistente y fuerza bruta; T1110 y T1021), 5715 y 5501
(acceso válido; T1078 y T1021), 5402 (\texttt{sudo} a root; T1548.003) y 5902
(cuenta creada; T1136, persistencia).

\subsection{Panel en vivo}\label{sec:panel}

\texttt{server.py} lanza un \texttt{tail -F} sobre
\texttt{/var/ossec/logs/alerts/alerts.json}, convierte cada línea en un objeto
reducido (regla, nivel, técnicas, IP de origen, usuario) y lo envía a los
navegadores conectados por \gls{sse} \parencite{whatwg-sse}. Además:

\begin{itemize}
  \item Guarda las últimas 300 alertas y las reenvía cuando un navegador se
        conecta, para que el panel no arranque vacío si se recarga.
  \item Manda un latido cada 15 segundos y le indica al navegador que se reconecte
        a los 2 segundos si se cae la conexión.
  \item Expone \texttt{/health} y \texttt{/api/recent} para comprobar el estado
        sin abrir el panel.
\end{itemize}

\texttt{index.html} mantiene el estado de la demo: el flujo de cuatro etapas, la
cuadrícula de las siete técnicas y el porcentaje. Una casilla pasa a
<<cerrado>> cuando llega una alerta real mapeada a esa técnica. El botón
<<Reiniciar demo>> pone todo a cero y guarda ese punto en el navegador.

\begin{nota}[Sobre la etapa <<Fable 5.1>> del panel]
El flujo del panel muestra cuatro etapas: Kali, Wazuh, Fable 5.1 y Regla. La
tercera es solo visual: repite el mapeo MITRE que ya trae la alerta de Wazuh. El
panel no llama a ningún modelo en tiempo real. Fable participó en la
construcción, no en la ejecución.
\end{nota}

\subsection{Scripts de apoyo}

\begin{itemize}
  \item \texttt{dashboard/run.sh}: arranca, detiene y revisa el panel
        (\texttt{start}, \texttt{stop}, \texttt{restart}, \texttt{status},
        \texttt{logs}).
  \item \texttt{scripts/close\_gap\_t1136.sh}: hace copia de seguridad de las
        reglas, añade la 100200 si no está y reinicia el manager.
  \item \texttt{scripts/enable\_active\_response.sh}: activa el bloqueo automático
        (\texttt{firewall-drop}) para alertas de nivel 10 o más
        \parencite{wazuh-ar}. Trae marcadores para poner las IPs de gestión en
        la lista blanca antes de ejecutarlo.
\end{itemize}

\section{La demo paso a paso}\label{sec:demo}

La demo empieza con el panel en cero y avanza por fases. La
\cref{tab:progresion} resume cómo subió la cobertura según las capturas.

\begin{table}[H]\centering\small
\begin{tabularx}{\textwidth}{L{3.4cm} C{1.9cm} C{1.6cm} Y}
\toprule
\filacab \thc{Fase} & \thc{Capturas} & \thc{Cobertura} & \thc{Técnicas nuevas y regla} \\
\midrule
Acceso válido y \texttt{sudo} & 02--04 & 43\,\% & T1078, T1021 (5715, 5501) y T1548, escalada de privilegios (5402) \\
Escaneo desde Kali & 05--08 & 71\,\% & T1595 y T1046 (100100, propia) \\
Intentos de acceso desde Kali & 09--13 & 86\,\% & T1110 (5710 y 5712) \\
Creación de cuenta en la EC2 & 14--18 & 100\,\% & T1136, persistencia (5902, de fábrica) \\
Regla 100200 aplicada & 20--22 & 100\,\% & T1136 detectada por la 100200 con nivel 10 \\
\bottomrule
\end{tabularx}
\caption{Progresión de la cobertura durante la demo.}\label{tab:progresion}
\end{table}

\evidenciaH[0.85\textwidth]{06_panel_71_escaneo_detectado.png}{El escaneo de Kali detectado por la
regla 100100: T1595 y T1046 pasan a <<cerrado>> y la cobertura sube al 71\,\%.}{fig:esc}

Los intentos de acceso por contraseña con una herramienta de fuerza bruta no
llegaron a probarse: el servidor rechaza la autenticación por contraseña. Por eso
esa fase se hizo con intentos de acceso de usuarios inexistentes, que Wazuh
registra con la regla 5710 y, al repetirse, con la 5712
(\cref{sec:problemas}).

La última fase corresponde a la persistencia mediante creación de cuentas
(T1136). La secuencia fue:

\begin{enumerate}
  \item Ejecuté \texttt{sudo useradd} en la EC2 para generar el evento de
        creación de cuenta (capturas 15 y 16).
  \item Apliqué la regla 100200, escrita por Fable, con
        \texttt{close\_gap\_t1136.sh} (captura 21).
  \item Al crear de nuevo una cuenta, la regla 100200 detectó el evento, lo
        mapeó a T1136 y lo elevó a nivel 10. En el log del panel aparece así:
        <<Persistencia: cuenta local creada (mapeada a ATT\&CK T1136 por
        Fable 5.1)>> (captura 22).
\end{enumerate}

El cierre del hueco consiste en eso: la creación de cuentas queda detectada y
etiquetada como la técnica de persistencia que es, con un nivel que la hace
visible frente al resto de eventos del sistema.

\evidenciaH[0.85\textwidth]{22_panel_100_regla_100200_t1136.png}{Panel al 100\,\% con la regla 100200
activa: la creación de \texttt{demo\_attacker2} entra como nivel 10, mapeada a T1136
(persistencia).}{fig:final}

\section{Decisiones de diseño}

\begin{decision}[Decisión 1 --- Python sin dependencias y SSE para el panel]
\contextodecision{El panel tenía que mostrar alertas en el momento en que ocurren
y arrancar en la misma EC2 sin instalar nada.}
\eleccion{Un servidor con solo la librería estándar de Python y \gls{sse} hacia el
navegador.}
\porque{El flujo va en un solo sentido (servidor a navegador), y SSE resuelve eso
sobre HTTP normal, con reconexión automática del navegador. Sin dependencias, el
panel se copia y se ejecuta en cualquier servidor con Python 3.}
\descartado{WebSockets (bidireccional, no hacía falta) y consultar la API del
indexer cada pocos segundos (añade retraso y credenciales).}
\end{decision}

\begin{decision}[Decisión 2 --- leer \texttt{alerts.json} con \texttt{sudo -n tail -F}]
\contextodecision{El archivo de alertas es de Wazuh y el panel no debía correr como
root.}
\eleccion{Lanzar \texttt{tail -F} con \texttt{sudo -n} desde un usuario normal.}
\porque{\texttt{tail -F} sigue al archivo aunque Wazuh lo rote cada día, y
\texttt{sudo -n} falla de inmediato en vez de quedarse esperando una contraseña.
Solo ese comando necesita privilegios, no todo el servidor web.}
\descartado{Ejecutar \texttt{server.py} como root.}
\end{decision}

\begin{decision}[Decisión 3 --- el manager se vigila a sí mismo]
\contextodecision{Hacía falta telemetría del servidor atacado.}
\eleccion{Usar los logs del propio host del manager, sin agente.}
\porque{El agente y el manager de Wazuh no se instalan en el mismo equipo, y el
manager ya analiza sus propios logs. Para una demo de un solo servidor era
suficiente.}
\descartado{Una segunda instancia con agente, que duplica costo y
complejidad.}
\end{decision}

\begin{decision}[Decisión 4 --- correlación en vez de una regla por evento]
\contextodecision{Un solo evento de conexión cerrada no significa escaneo;
pasa todo el tiempo.}
\eleccion{Eventos base de nivel 1 y una regla de correlación (100100) que exige dos
o más desde la misma IP en 5 minutos.}
\porque{Así cada conexión suelta no genera ruido, pero un escaneo sí se alerta, y
con la IP del escáner.}
\descartado{Alertar cada conexión cerrada, que llenaría el panel de falsos
positivos.}
\end{decision}

\begin{decision}[Decisión 5 --- la regla 100200 en un archivo aparte]
\contextodecision{La demo necesitaba un momento visible de <<regla nueva
aplicada>>.}
\eleccion{Dejar la 100200 fuera de \texttt{local\_rules.xml} y aplicarla en vivo con
\texttt{close\_gap\_t1136.sh}.}
\porque{Así se puede mostrar el antes y el después sin editar archivos en el
escenario. El script hace copia de seguridad y no duplica la regla si ya está.}
\descartado{Editar las reglas a mano durante la presentación.}
\end{decision}

\section{Problemas encontrados}\label{sec:problemas}

\begin{problema}[Problema 1 --- Wazuh de fábrica no detectaba el escaneo]
\sintoma{El escaneo desde Kali funcionó: llegó al servidor, encontró los puertos
22 y 443 abiertos e identificó los servicios (captura 05). Aun así, el panel no
mostraba ninguna alerta para T1595 ni T1046.}
\hipotesis{Los eventos del escaneo sí llegaban a Wazuh, pero ninguna regla de
fábrica los interpretaba como escaneo.}
\verificacion{La regla de fábrica para conexión cerrada (5722) tiene nivel 0. Las
sondas no-SSH aparecían en el log como \texttt{banner exchange \ldots\ invalid
format}, un formato que ningún decoder extraía.}
\resultado{Fable generó el decoder \texttt{sshd-banner-invalid} y las reglas
100098, 100099 y 100100, y las ajustó contra los eventos reales del servidor. El
escaneo, que antes era invisible para Wazuh, ahora produce una alerta de nivel 8
con T1046 y T1595, y la cobertura sube del 43\,\% al 71\,\% (capturas 06 a 08).}
\end{problema}

\begin{problema}[Problema 2 --- la fuerza bruta por contraseña no era posible]
\sintoma{La herramienta de fuerza bruta respondía que el servidor no admite
autenticación por contraseña (captura 10).}
\hipotesis{Es el endurecimiento del propio servidor: SSH solo con clave.}
\verificacion{El mensaje de error lo confirma. Relajar esa configuración para la
demo no tenía sentido.}
\resultado{La fase se hizo con intentos de acceso de usuarios inexistentes, que
Wazuh sí registra (5710) y correlaciona como fuerza bruta (5712, T1110).}
\end{problema}

\begin{problema}[Problema 3 --- T1136: la hipótesis inicial era incorrecta]
\sintoma{Antes de crear cualquier cuenta, la casilla T1136 estaba en rojo.}
\hipotesis{Se asumió que Wazuh detectaba la creación de cuentas (regla 5902) pero no
la mapeaba a MITRE, y con esa idea se escribió la regla 100200 para
<<cerrar el hueco>>.}
\verificacion{En la captura 14, la alerta de la 5902 ya trae la etiqueta T1136 y la
casilla pasa a verde antes de aplicar la 100200. El código fuente de Wazuh 4.14
lo confirma: la regla 5902 incluye \texttt{<mitre><id>T1136</id></mitre>}
\parencite{wazuh-syslog-rules}. La casilla estaba en rojo porque aún no se había
creado ninguna cuenta, no porque faltara el mapeo.}
\resultado{La 100200 sigue siendo útil, pero por otra razón: con ella activa, la
creación de la cuenta queda registrada como <<Persistencia: cuenta local creada
(mapeada a ATT\&CK T1136 por Fable 5.1)>>, con nivel 10 en vez de 8, lo que la pone
dentro del umbral de respuesta activa (captura 22). La lección es que un <<hueco>> en la cuadrícula significa
<<sin alerta todavía>>, y eso hay que comprobarlo contra el ruleset antes de
escribir una regla.}
\end{problema}

\begin{problema}[Problema 4 --- fugas en capturas ya censuradas]
\sintoma{Las capturas se censuraron con un script (bloque opaco y reemplazo por
direcciones RFC 5737), y una verificación con OCR dio cero coincidencias.}
\hipotesis{El OCR no ve todo.}
\verificacion{Revisadas a ojo y con el contraste subido, tres capturas todavía
mostraban datos reales: una terminal transparente dejaba ver detrás un comando
anterior con la IP del servidor, y en otras dos la última fila del log, cortada
en el borde, mostraba parte de una IP y de un nombre de host internos.}
\resultado{Las copias del repositorio se recortaron y se revisaron de nuevo. La
lección quedó en el checklist OPSEC de mi plantilla.}
\end{problema}

\section{Seguridad del propio proyecto}\label{sec:seguridad}

\begin{riesgo}[El panel no tiene autenticación]
\texttt{server.py} escucha en \texttt{0.0.0.0:8080} y muestra alertas reales,
incluidas IPs de origen y usuarios. El grupo de seguridad de la EC2 debe permitir
el puerto 8080 solo desde la IP de quien presenta.
\end{riesgo}

\begin{riesgo}[El bloqueo automático puede bloquear al administrador]
\texttt{enable\_active\_response.sh} bloquea por firewall las IPs que disparan
alertas de nivel 10 o más. Antes de ejecutarlo hay que reemplazar los marcadores
\texttt{IP\_GESTION\_SSH} e \texttt{IP\_GESTION\_NAVEGADOR}, o se pierde el acceso
al servidor.
\end{riesgo}

Para publicar el proyecto:

\begin{itemize}
  \item En el código y la documentación, la IP del servidor aparece como
        \texttt{TU\_IP\_EC2}. No hay claves, contraseñas ni IPs de gestión.
  \item Las capturas usan direcciones de documentación \parencite{rfc5737}:
        \texttt{198.51.100.20} para el servidor, \texttt{203.0.113.10} para el
        atacante y \texttt{198.51.100.30} para el origen de las sesiones de
        administración. El nombre del host aparece como \texttt{wazuh-lab}.
  \item El \texttt{.gitignore} excluye claves, archivos \texttt{.env},
        credenciales, logs de alertas y notas internas de OPSEC.
\end{itemize}

Además, la instancia sigue generando costo apagada (disco EBS e IP pública), así
que al terminar conviene liberar la IP y terminar la instancia.

\section{Limitaciones y trabajo futuro}

\begin{limitacion}[El modelo no participa en la ejecución]
La etapa <<Fable 5.1>> del panel es visual. Un siguiente paso real sería que el
modelo recibiera las alertas sin técnica y propusiera la regla en vivo, con
aprobación humana antes de aplicarla.
\end{limitacion}

\begin{limitacion}[Alcance acotado]
Un solo servidor, siete técnicas y ataques controlados. La cobertura mide
<<técnica con al menos una alerta real>>, no la calidad de la detección ante
variantes del mismo ataque. Además, la \texttt{t3.large} (2 vCPU) está por debajo
de los 4 vCPU que Wazuh recomienda para el modo todo en uno
\parencite{wazuh-quickstart}. Alcanzó para la demo, pero no sirve como referencia
de rendimiento.
\end{limitacion}

\begin{recomendacion}[Modo de anonimización en el panel]
Añadir en \texttt{server.py}, donde se arma cada alerta, un modo activado con
\texttt{ANON=1} que reemplace IPs y nombres de host antes de enviarlos por SSE.
Así se puede grabar o proyectar el panel sin censurar después.
\end{recomendacion}

\section{Conclusiones}

\begin{enumerate}
  \item \textbf{Un modelo de IA puede escribir y mejorar reglas de detección de
        forma autónoma.} Es el aporte principal del proyecto. Dentro del flujo
        que diseñé (SOC montado, ataques definidos y validación contra alertas
        reales), Fable 5.1 generó el decoder y las reglas que convirtieron un
        escaneo invisible para Wazuh en una alerta mapeada a MITRE, y las ajustó
        hasta que funcionaron sobre los eventos reales. Mi papel fue diseñar el
        flujo, dirigir y verificar; el de Fable, producir la detección.
  \item \textbf{El hueco real estaba en la detección del escaneo.} Las reglas de fábrica cubren
        bien los accesos, la escalada de privilegios con \texttt{sudo} (T1548) y
        la persistencia por creación de cuentas (T1136), pero no
        correlacionan un escaneo de red contra OpenSSH 9.x. Ahí fue donde la
        detección como código aportó.
  \item \textbf{Un panel vacío no es un hueco.} El caso T1136 mostró que hay que
        verificar contra el ruleset antes de dar por hecho que algo no se
        detecta.
  \item \textbf{La autonomía funciona mejor con un flujo que la verifique.} Que
        cada regla se probara contra un Wazuh real, y que la evidencia se revisara
        hasta encontrar el error de la hipótesis sobre T1136, fue lo que hizo
        confiable el resultado.
  \item \textbf{Publicar también es parte del proyecto.} La revisión OPSEC encontró
        fugas después de una verificación automática que decía que no las había.
\end{enumerate}

%%%%%%%%%%%%%%%%%%%%%%%%%%%%%%%%%%%%%%%%%%%%%%%%%%%%%%%%%%%%%%%%%%%%%%%%%%%%%%%
%% CAPA 3 — ANEXOS
%%%%%%%%%%%%%%%%%%%%%%%%%%%%%%%%%%%%%%%%%%%%%%%%%%%%%%%%%%%%%%%%%%%%%%%%%%%%%%%
\clearpage
\iniciaranexos

\section{Matriz de cobertura ATT\&CK}\label{anx:matriz}

\begin{table}[H]\centering\small
\begin{tabularx}{\textwidth}{L{3.1cm} L{2.4cm} Y C{1.6cm} C{1.2cm} C{1.5cm}}
\toprule
\filacab \thc{Técnica} & \thc{Táctica} & \thc{Actividad en la demo} & \thc{Regla} & \thc{Captura} & \thc{Estado} \\
\midrule
T1595 Active Scanning & Reconocimiento & Escaneo desde Kali & 100100 & 06 & \hecho \\
T1046 Network Service Discovery & Descubrimiento & Escaneo desde Kali & 100100 & 06 & \hecho \\
T1110 Brute Force & Acceso a credenciales & Accesos con usuarios inexistentes desde Kali & 5710, 5712 & 12, 13 & \hecho \\
T1078 Valid Accounts & Acceso inicial & Acceso SSH válido & 5715, 5501 & 03, 04 & \hecho \\
T1021 Remote Services & Movimiento lateral & Acceso SSH & 5715, 5710 & 03, 12 & \hecho \\
T1548 Abuse Elevation Control & Escalada de privilegios & \texttt{sudo} a root & 5402 & 04 & \hecho \\
T1136 Create Account & Persistencia & Creación de cuenta local con \texttt{useradd} en la EC2 & 5902, 100200 & 14, 22 & \hecho \\
\bottomrule
\end{tabularx}
\caption{Cobertura de las siete técnicas del set. La táctica es la principal en el
contexto de la demo; T1078 figura en ATT\&CK bajo varias tácticas.}\label{tab:matriz}
\end{table}

\section{Reglas y decoder}\label{anx:reglas}

\textbf{Decoder} (\texttt{wazuh-rules/local\_decoder.xml}):

\begin{lstlisting}[language=XML]
<decoder name="sshd-banner-invalid">
  <parent>sshd</parent>
  <prematch>banner exchange: Connection from </prematch>
  <regex offset="after_prematch">^(\S+) port (\d+):</regex>
  <order>srcip, srcport</order>
</decoder>
\end{lstlisting}

\textbf{Reglas de escaneo} (\texttt{wazuh-rules/local\_rules.xml}):

\begin{lstlisting}[language=XML]
<rule id="100099" level="1">
  <if_sid>5722</if_sid>
  <description>sshd: conexion cerrada sin autenticar (evento base).</description>
  <group>recon_base,</group>
</rule>
<rule id="100098" level="1">
  <if_sid>5700</if_sid>
  <match>banner exchange: Connection from</match>
  <description>sshd: sonda con protocolo invalido (evento base).</description>
  <group>recon_base,</group>
</rule>
<rule id="100100" level="8" frequency="2" timeframe="300">
  <if_matched_group>recon_base</if_matched_group>
  <same_srcip />
  <description>Recon: posible escaneo contra SSH desde $(srcip).</description>
  <mitre><id>T1046</id><id>T1595</id></mitre>
  <group>recon,network_scan,</group>
</rule>
\end{lstlisting}

\textbf{Regla de creación de cuentas} (\texttt{wazuh-rules/local\_rules\_t1136.xml}):

\begin{lstlisting}[language=XML]
<rule id="100200" level="10">
  <if_sid>5902</if_sid>
  <description>Persistencia: cuenta local creada (mapeada a ATT&amp;CK T1136 por Fable 5.1).</description>
  <mitre><id>T1136</id></mitre>
  <group>account_manipulation,persistence,</group>
</rule>
\end{lstlisting}

Las descripciones están abreviadas; el texto completo está en los archivos del
repositorio.

\section{Montaje desde cero}\label{anx:montaje}

\textbf{1. Instancia.} EC2 con Amazon Linux 2023 e IP fija. En el grupo de
seguridad, abrir 22 y 443 solo desde tu IP, y 8080 solo desde la IP del navegador
del panel.

\textbf{2. Wazuh todo en uno} \parencite{wazuh-quickstart}:

\begin{lstlisting}[style=terminal]
curl -sO https://packages.wazuh.com/4.14/wazuh-install.sh
sudo bash ./wazuh-install.sh -a
\end{lstlisting}

\textbf{3. Reglas y decoder propios:}

\begin{lstlisting}[style=terminal]
sudo cp wazuh-rules/local_rules.xml   /var/ossec/etc/rules/local_rules.xml
sudo cp wazuh-rules/local_decoder.xml /var/ossec/etc/decoders/local_decoder.xml
sudo chown wazuh:wazuh /var/ossec/etc/rules/local_rules.xml \
                       /var/ossec/etc/decoders/local_decoder.xml
sudo systemctl restart wazuh-manager
\end{lstlisting}

\textbf{4. Panel.} El usuario que lo arranca debe poder ejecutar
\texttt{sudo -n tail} sobre \texttt{alerts.json}.

\begin{lstlisting}[style=terminal]
cd dashboard
./run.sh start     # http://TU_IP_EC2:8080
./run.sh status    # consulta /health
\end{lstlisting}

\textbf{5. Regla 100200 (opcional, en vivo):}
\texttt{./scripts/close\_gap\_t1136.sh}.

\section{Evidencia}\label{anx:evidencia}

Las 22 capturas están en \texttt{evidencia/}, numeradas en el orden en que
ocurrieron. Aquí van las que sostienen los resultados, en el orden de la demo;
todas usan direcciones de documentación.

\Needspace{22\baselineskip}\textbf{Punto de partida (43\,\%).} El panel ya muestra cerradas T1078, T1021 y
T1548. Las produjeron un acceso SSH válido y un \texttt{sudo} a root, y las
detectaron reglas de fábrica (5715, 5501 y 5402). Es la base que Wazuh cubre solo,
antes de cualquier regla nueva.

\evidenciaH{03_panel_43_con_log.png}{Captura 03. Primeras técnicas cerradas (43\,\%)
con accesos válidos y \texttt{sudo}.}{fig:ev03}

\Needspace{22\baselineskip}\textbf{El escaneo llega al servidor.} Desde Kali se lanzó \texttt{nmap} contra la
EC2: encontró los puertos 22 y 443 abiertos e identificó OpenSSH 9.9 y el servicio
web de Wazuh. El ataque funcionó; lo que faltaba era que Wazuh lo detectara, y eso
es lo que resuelven las reglas de Fable (Resultado 2).

\evidenciaH[0.8\textwidth]{05_kali_nmap_escaneo.png}{Captura 05. Escaneo desde Kali
contra el servidor (huellas SSH tapadas).}{fig:ev05}

\Needspace{22\baselineskip}\textbf{El endurecimiento hace su trabajo.} El intento de fuerza bruta por contraseña
termina con el servidor respondiendo que no admite ese método de autenticación. Por
eso la fase de fuerza bruta se hizo con usuarios inexistentes (Problema 2).

\evidenciaH[0.8\textwidth]{10_kali_hydra_rechazado_ssh_solo_clave.png}{Captura 10.
El servidor rechaza la autenticación por contraseña.}{fig:ev10}

\Needspace{22\baselineskip}\textbf{Fuerza bruta detectada (86\,\%).} Los intentos de acceso con usuarios
inexistentes desde Kali aparecen en el log como regla 5710, mapeados a T1110 y
T1021, y T1110 pasa a cerrada. Para entonces el escaneo ya figura como detectado
por la regla 100100.

\evidenciaH{12_panel_86_fuerza_bruta_detectada.png}{Captura 12. Intentos con
usuarios inexistentes detectados por la 5710 (86\,\%).}{fig:ev12}

\Needspace{22\baselineskip}\textbf{Persistencia por creación de cuenta (100\,\%).} El \texttt{sudo useradd}
ejecutado en la EC2 genera la alerta 5902 con T1136, y el panel llega al 100\,\%.
Esta captura es la evidencia del Problema 3: la regla de fábrica ya etiquetaba
T1136 antes de aplicar la 100200.

\evidenciaH{14_panel_100_cuenta_creada_regla_5902.png}{Captura 14. La regla de fábrica
5902 ya trae T1136: la casilla se cierra antes de aplicar la 100200.}{fig:ev14}

\Needspace{22\baselineskip}\textbf{Despliegue de la regla de Fable.} En la EC2 se ejecuta
\texttt{close\_gap\_t1136.sh}, que agrega la regla 100200 y reinicia el manager. A
partir de aquí, cada cuenta nueva se registra como <<Persistencia: cuenta local
creada (mapeada a ATT\&CK T1136 por Fable 5.1)>> con nivel 10, como se ve en la
captura 22 (\cref{fig:final}).

\evidenciaH[0.75\textwidth]{21_ec2_close_gap_t1136.png}{Captura 21. Aplicación de la
regla 100200 con \texttt{close\_gap\_t1136.sh}.}{fig:ev21}

\clearpage
\printglossary[type=\acronymtype,title={Glosario de siglas},toctitle={Glosario de siglas}]
\printbibliography[title={Referencias},heading=bibintoc]
\end{document}
